\clearpage
\fancyhead[L]{\small GAUSS AN BODE.}
\hypertarget{native-298-lower}{}
\begin{center}
GAUSS AN BODE.
\end{center}

\begin{center}
\begin{tabular}{c}
\hline
Astronomisches Jahrbuch für 1811. S. 135--139. Berlin 1808.\\
\hline
\end{tabular}
\end{center}

\begin{flushright}
Göttingen 1808. Januar 24.
\end{flushright}

{\rightskip=0pt plus 1em\par\indent Anfangs die mit der Veränderung meines Aufenthalts; verbundenen Unruhen, und nachher die
durch eine höchst ungünstige Witterung vereitelte Hoffnung, Ihre communicirten
Cometenbeobachtungen durch zahlreiche hiesige erwidern zu können, werden mein langes
Stillschweigen entschuldigen. Hier haben Ew. --- indess die wenigen Beobachtungen des
Cometen, die ich dem hier fast ununterbrochen bedeckten Himmel habe abgewinnen können.\par}

\begin{center}\fontsize{9}{11}\selectfont\setlength{\tabcolsep}{2pt}
\begin{tabular}{lrrrl}
& M. Z. in Göttingen. & AR. & Decl. & \\
1807. Dec. 17. & $11^{\text{u}}\,24^{\text{m}}\,44^{\text{s}}$ & $305^\circ\;54'\;0''$ & $45^\circ\;4'\;30''$ & \\
29. & $7\;59\;21$ & $318\;57\;41$ & $46\;41\;48$ & \\
& $11\;36\;33$ & $319\;6\;35$ & $46\;43\;11$ & \\
1808. Jan. 4. & $7\;22\;12$ & $325\;11\;19$ & $47\;12\;46$ & \smash{\raisebox{16.5pt}{$\left.\vphantom{\begin{array}{c}0\\0\\0\\0\end{array}}\right\}$ von mir.}}\\
& $9\;13\;57$ & $325\;15\;53$ & & \\
& $17\;30\;26$ & $325\;36\;38$ & & von Hrn. Prof. \textsc{Harding}.\\
\end{tabular}
\end{center}

\par\indent Die parabolische Bahn des Cometen hatte ich bereits im Anfang November aus einigen Bremer
und Lilienthaler Beobachtungen berechnet; da dieselben bisher noch nicht bekannt geworden
sind, so theile ich sie Ihnen hier mit.

\SourceLeaf{299}\noindent Durchgang durch das Perihelium 18. Sept. $19^{\text{u}}\,16^{\text{m}}\,0^{\text{s}}$ Merid. von Paris.


\begin{center}\fontsize{10}{12}\selectfont\setlength{\tabcolsep}{3pt}\renewcommand{\arraystretch}{1.15}\begin{tabular}{lr}
Länge des Perihelium .\;.\;.\;.\;.\;.\;.\;.\;.\;.\;.\;.\;.\; & $271^\circ\;0'\;13''3.$ \\
Logarithm des kleinsten Abstandes , .\;.\;.\;.\;.\;.\;.\;.\;.\;.\; & $9.8114927$ \\
Aufsteigender Knoten .\;.\;.\;.\;.\;.\;.\;.\;.\;.\;.\;.\;.\; & $266^\circ\;38'\;31''2.$ \\
Neigung der Bahn .\;.\;.\;.\;.\;.\;.\;.\;.\;.\;.\;.\;.\;.\; & $63\;12\;35.6.$ \\
\end{tabular}\end{center}

\noindent Bewegung rechtläufig.


\par\indent Diese Elemente wichen Ende November und Anfang December gegen $1'$ von den Beobachtungen ab,
nachher aber hat der Fehler wieder abgenommen. Sollte es mir gelingen noch einige
Beobachtungen des Cometen zu erhalten, so werde ich danach obige Bahn noch einmal verbessern,
die Correctionen werden aber gewiss sehr klein sein. Ueber die Abweichung der Bahn von der
Parabel, die nicht beträchtlich sein kann, wird sich aus den Beobachtungen schwerlich etwas
entscheiden lassen.

\par\indent Meine neuen (X.) Elemente der \emph{Pallas} sind folgende:


\begin{center}\fontsize{10}{12}\selectfont\setlength{\tabcolsep}{3pt}\renewcommand{\arraystretch}{1.15}\begin{tabular}{lr}
Epoche 1808 Meridian von Göttingen .\;.\;.\;.\;.\;.\;.\; & $252^\circ\;32'\;28''5$ \\
Tägliche tropische Bewegung .\;.\;.\;.\;.\;.\;.\;.\;.\;.\;.\;.\; & $770''2143$ \\
Logarithm der halben grossen Axe .\;.\;.\;.\;.\;.\;.\;.\;.\; & $0.4423149$ \\
Excentricität .\;.\;.\;.\;.\;.\;.\;.\;.\;.\;.\;.\;.\;.\;.\;.\;.\; & $0.2450198$ \\
Perihelium 1803 .\;.\;.\;.\;.\;.\;.\;.\;.\;.\;.\;.\;.\;.\;.\; & $121^\circ\;3'\;11''4$ \\
$\AscendingNode\;1803$ .\;.\;.\;.\;.\;.\;.\;.\;.\;.\;.\;.\;.\;.\;.\;.\;.\;.\; & $172\;28\;56.9$ \\
Neigung der Bahn .\;.\;.\;.\;.\;.\;.\;.\;.\;.\;.\;.\;.\;.\; & $34\;37\;41.0$ \\
\end{tabular}\end{center}


\par\indent Die Elemente der \emph{Juno} habe ich vor Kurzem nach Hrn. \textsc{Bessel}'s im vorigen Jahre
angestellten Beobachtungen verbessert. Hier die neuen (VII.) Elemente:


\begin{center}\fontsize{10}{12}\selectfont\setlength{\tabcolsep}{3pt}\renewcommand{\arraystretch}{1.15}\begin{tabular}{lr}
Epoche, Meridian von Göttingen 1805 .\;.\;.\;.\;.\;.\;.\; & $42^\circ\;37'\;3''7$ \\
1806 .\;.\;.\;.\;.\;.\;.\; & $125\;7\;57.8$ \\
1807 .\;.\;.\;.\;.\;.\;.\; & $207\;38\;51.8$ \\
1808 .\;.\;.\;.\;.\;.\;.\; & $290\;23\;19.8$ \\
Mittlere tägliche tropische Bewegung .\;.\;.\;.\;.\;.\;.\;.\;.\; & $813''8468$ \\
Tropische Umlaufszeit .\;.\;.\;.\;.\;.\;.\;.\;.\;.\;.\; & $1592\;\text{Tage}\;10\tfrac12\;\text{St.}$ \\
Sonnennähe 1805 .\;.\;.\;.\;.\;.\;.\;.\;.\;.\;.\;.\;.\;.\;.\; & $53^\circ\;19'\;0''2$ \\
Aufsteigender Knoten 1805 .\;.\;.\;.\;.\;.\;.\;.\;.\;.\;.\; & $171\;4\;28.2$ \\
Excentricität .\;.\;.\;.\;.\;.\;.\;.\;.\;.\;.\;.\;.\;.\;.\;.\;.\; & $0.2554996$ \\
Logarithm der halben grossen Axe .\;.\;.\;.\;.\;.\;.\;.\;.\; & $0.4263781$ \\
Neigung der Bahn .\;.\;.\;.\;.\;.\;.\;.\;.\;.\;.\;.\;.\;.\; & $13^\circ\;4'\;26.2$ \\
\end{tabular}\end{center}


\par\indent Die hienach berechnete Ephemeride für 1808 ist folgende: [Sie gibt für Mitternacht in
Göttingen vom 16. April bis 28. December mit vier Tagen Zwischenräumen AR. und Decl. auf
Minuten genau.]

\par\indent Mein Werk über die Bestimmung der Bahnen der Himmelskörper ist jetzt unter der Presse; auf
Ostern hoffe ich wird es erscheinen können.

\SourceLeaf{300}
\begin{flushright}
Göttingen 1808 Juli 23.
\end{flushright}

\par\indent Ew. --- Wunsche zufolge theile ich Ihnen hier meine Beobachtungen mit, die ich seither an
den neuen Planeten habe machen können. Sie sind alle am Kreismicrometer gemacht; an unserm
Mauerquadranten wird dies Jahr blos die Vesta beobachtet werden können, und diese culminirt
noch bei Tage.

\Needspace{9\baselineskip}
\begin{center}\emph{Beobachtungen der Juno.}\end{center}

\begin{center}\fontsize{10}{12}\selectfont\setlength{\tabcolsep}{3pt}\renewcommand{\arraystretch}{1.15}\begin{tabular}{lrrr}
1808 & \shortstack{Mittl. Zeit\\in Göttingen} & \shortstack{Gerade\\Aufsteigung} & \shortstack{südliche\\Abweichung}\\
Juni 20 & $11^{\text{u}}\,49^{\text{m}}\,0^{\text{s}}$ & $315^\circ\;29'\;34''$ & $2^\circ\;16'\;23''$\\
22 & $12\;0\;45$ & $315\;23\;1$ & $2\;14\;29$\\
Juli 6 & $12\;42\;23$ & $314\;1\;0$ & $2\;19\;7$\\
\end{tabular}\end{center}

\Needspace{9\baselineskip}\begin{center}\emph{Beobachtungen der Vesta.}\end{center}

\begin{center}\fontsize{9}{12}\selectfont\setlength{\tabcolsep}{3pt}\renewcommand{\arraystretch}{1.15}\begin{tabular}{lrrrl}
Juni 22 & $13^{\text{u}}\,46^{\text{m}}\,25^{\text{s}}$ & $353^\circ\;12'\;17''$ & $9^\circ\;12'\;16''$ & \\
Juli 1 & $14\;38\;57$ & $354\;28\;3$ & $9\;9\;39$ & \\
2 & $12\;50\;2$ & $355\;7\;7$ & & (v. Hrn. Prof. \textsc{Harding}).\\
13 & $13\;22\;37$ & $356\;35\;27$ & $9\;27\;40$ & \\
\end{tabular}\end{center}

\Needspace{7\baselineskip}\begin{center}\emph{Beobachtungen der Pallas.}\end{center}

\begin{center}\fontsize{9}{12}\selectfont\setlength{\tabcolsep}{3pt}\renewcommand{\arraystretch}{1.15}\begin{tabular}{lrrrl}
Juli 13 & $12^{\text{u}}\,3^{\text{m}}\,17^{\text{s}}$ & $300^\circ\;21'\;49''$ & $18^\circ\;41'\;44''$ Nordl. & \\
15 & $12\;1\;28$ & $299\;57\;3$ & $18\;32\;40$ & (v. Hrn. Prof. \textsc{Harding}).\\
\end{tabular}\end{center}

\par\indent Die Vesta hat jetzt schon die $7^{\text{te}}$ Grösse; aber Juno und Pallas sind sehr
lichtschwach und kaum $10^{\text{ter}}$ Grösse zu schätzen. Ceres habe ich noch nicht
gesehen. Die Beobachtungen der Pallas werden nicht sehr zuverlässig sein, da sie nur mit
einem schwächern Instrumente haben gemacht werden können, indem der Planet für das
Herschelsche Telescop jetzt zu hoch steht. --- Wegen der Elemente beziehe ich mich auf
No. 14. 40. 107. unsrer gel. Anzeigen [Seite 292. 299. 300 dieses Bandes]. Ephemeriden für
das nächste Jahr werde ich erst nach neuer Verbesserung der Elemente besorgen. Bei der
Vesta finde ich um so weniger eine nöthig, da sie, wenn sie im März 1809 unsichtbar
geworden, erst im Oct. oder Nov. 1809 wieder sichtbar werden wird. --- Hier noch einige
Jupitertrabanten-Eintritte:

\begin{center}
\fontsize{10}{12}\selectfont\setlength{\tabcolsep}{3pt}\begin{tabular}{lccl}
$2^{\text{ter}}$ Trab. Juni 22. & $12^{\text{u}}\,45^{\text{m}}\,9^{\text{s}}$ & & \textsc{Harding} $3\tfrac12$f. \textsc{Dollond}.\\
& $12\;45\;32$ & & \textsc{Gauss} $10$f. \textsc{Herschel}.\\
$1^{\text{ster}}$ Trab. Juli 13. & $13\;45\;36$ & & \textsc{Harding} $4$f. \textsc{Dollond}.\\
& $13\;45\;38$ & & \textsc{Thiarks} $3\tfrac12$f. \textsc{Dollond}.\\
& $13\;45\;38$ & & \textsc{Gauss} $10$f. \textsc{Herschel}.\\
\end{tabular}

Austritte.

\fontsize{10}{12}\selectfont\setlength{\tabcolsep}{3pt}\begin{tabular}{lccl}
$4^{\text{ter}}$ Trab. Juli 14. & $13\;35\;16$ & & \textsc{Harding} $10$f. \textsc{Herschel}.\\
& $13\;35\;56$ & & \textsc{Thiarks} $3\tfrac12$f. \textsc{Dollond}.\\
$2^{\text{ter}}$ Trab. Juli 24. & $12\;32\;13$ & & \textsc{Thiarks} $3\tfrac14$f. \textsc{Dollond}.\\
& $12\;32\;27$ & & \textsc{Harding} $4$f. \textsc{Dollond}.\\
\end{tabular}
\end{center}

\Needspace{8\baselineskip}\par\indent Ferner 2 Sternbedeckungen.

\smallskip
\noindent Juni 4.\quad $9^{\text{u}}\,3^{\text{m}}\,39^{\text{s}}0$\quad Eintritt von $\iota$ Virginis. \textsc{Gauss} und \textsc{Harding} in derselben Secunde.

\smallskip
\noindent Juli 6.\quad $10\;39\;6.2$\quad Eintritt von $\mu$ Sagittarii. \textsc{Gauss}.

\smallskip
{\rightskip=0pt plus 3em\par\indent Kürzlich ist der Instrumentenvorrath unsrer Sternwarte durch ein Geschenk des Königs mit
einem Spiegeltelescop von \textsc{Chevalier} und einem \textsc{Berthoud}'schen Zeithalter vermehrt. --- Der
Druck meines
\SourceLeaf{301}Werks über die Bestimmung der Planetenbahnen ist noch nicht vollendet. Erst 12 Bogen sind
bis jetzt fertig, das Ganze wird etwa 30 halten. Von Hrn. Prof. \textsc{Hardings} Karten erwarten
wir posttäglich das $4^{\text{te}}$ Blatt. Möchte nur dieses kostbare Unternehmen auch eine
hinlängliche Unterstützung bei dem Publicum finden! ---\par}
